

Comme @@M0@@ est à graduation finie, les gradués associés aux filtrations sur @@M1@@ de @@M2@@-anneau et par le degré sont identiques (6.6.7). On définit alors @@M3@@ par l'identité sur @@M4@@ et par :
@@M5@@
et on prolonge par additivité. D'après la relation (6.5.3), c'est un homomorphisme d'anneaux gradués, et il est facile de voir qu'il est compatible aux opérations de @@M6@@. Le morphisme inverse @@M7@@ peut être défini ainsi : @@M8@@, @@M9@@, le @@M10@@-ième coefficient @@M11@@ de @@M12@@ ne dépend pas de la façon dont @@M13@@ a été remonté dans @@M14@@, et on pose :
@@M15@@
L'homomorphisme @@M16@@ est donc un isomorphisme.

L'isomorphisme @@M17@@ est défini par la classe de Chern complétée @@M18@@. Pour vérifier qu'elle définit bien un isomorphisme, il suffit de vérifier que l'homomorphisme obtenu par passage aux gradués associés en est un, les filtrations étant discrètes. On obtient donc :
@@M19@@
Il suffit de montrer que le composé est l'identité. Comme les homomorphismes sont des homomorphismes de @@M20@@-algèbres graduées, il suffit, d'après la définition de la filtration de @@M21@@-anneau de @@M22@@, de montrer que c'est l'identité pour les @@M23@@, avec @@M24@@. On obtient alors :
@@M25@@
Or @@M26@@, d'après (6.6.4); donc :
@@M27@@
